%=====================================================================
% Modelo de datos · Llaves foráneas
%=====================================================================
\subsection*{Llaves foráneas}
\addcontentsline{toc}{subsection}{Llaves foráneas}

La columna \emph{Obligatoria} indica si la llave foránea admite nulo. Ninguna llave foránea borra ni actualiza en cascada: la fila de \textit{Usuario} nunca se elimina (\mbox{D-07}), y las pocas filas que se eliminan físicamente no tienen filas que dependan de ellas.

Algunas llaves foráneas son compuestas porque, además de exigir que la fila referida exista, hacen cumplir una regla de negocio:

\begin{reglas}
  \item \textbf{Solo se equipa lo que se posee.} \textit{Equipamiento} apunta a \textit{UsuarioObjeto} con (\at{id\_usuario}, \at{id\_objeto}), no a \textit{ObjetoCosmetico} (\mbox{CU-14} \mbox{RN-02}).
  \item \textbf{El objeto es de su ranura.} \textit{Equipamiento} y \textit{ParticipacionObjeto} apuntan a \textit{ObjetoCosmetico} con (\at{id\_objeto}, \at{id\_ranura}), así que no puede equiparse un marco en la ranura del peón.
  \item \textbf{Solo los participantes actúan en la partida.} \textit{Jugada}, \textit{OfertaTablas} y \textit{EnlaceEspectador} apuntan a \textit{Participacion} con (\at{id\_partida}, \at{id\_usuario}): solo quien juega la partida mueve, ofrece tablas o comparte el enlace. En \textit{Jugada}, un \at{id\_usuario} nulo, que es la IA, no se comprueba.
  \item \textbf{Las pruebas de un reporte son de una partida que jugaron los dos.} \textit{Reporte} apunta dos veces a \textit{Participacion}: con (\at{id\_partida}, \at{id\_denunciante}) y con (\at{id\_partida}, \at{id\_reportado}). Si se adjunta una partida, el denunciante y el reportado tuvieron que jugarla (\mbox{CU-42} \mbox{RN-09}); si no se adjunta, \at{id\_partida} es nula y ninguna de las dos se comprueba.
  \item \textbf{El historial de división pertenece a unas estadísticas.} \textit{HistorialDivision} apunta a \textit{EstadisticaModo} con (\at{id\_usuario}, \at{id\_modo}), no por separado a \textit{Usuario} y a \textit{Modo}.
\end{reglas}

% Archivo generado por bd/generar_tablas.ps1 a partir de bd/bastion.sql. No editar a mano.
\begin{center}\small
\begin{longtable}{|L{0.21\textwidth}|L{0.26\textwidth}|L{0.29\textwidth}|L{0.12\textwidth}|}
\hline
\textbf{Entidad} & \textbf{Llave foránea} & \textbf{Referencia} & \textbf{Obligatoria} \\ \hline
\endfirsthead
\hline
\textbf{Entidad} & \textbf{Llave foránea} & \textbf{Referencia} & \textbf{Obligatoria} \\ \hline
\endhead
\textit{Objeto\-Cosmetico} & (\texttt{id\_\allowbreak{}ranura}) & \textit{Ranura} (\texttt{id\_\allowbreak{}ranura}) & Sí \\ \hline
\textit{Tipo\-Caja\-Objeto} & (\texttt{id\_\allowbreak{}tipo\_\allowbreak{}caja}) & \textit{Tipo\-Caja} (\texttt{id\_\allowbreak{}tipo\_\allowbreak{}caja}) & Sí \\ \hline
\textit{Tipo\-Caja\-Objeto} & (\texttt{id\_\allowbreak{}objeto}) & \textit{Objeto\-Cosmetico} (\texttt{id\_\allowbreak{}objeto}) & Sí \\ \hline
\textit{Sesion} & (\texttt{id\_\allowbreak{}usuario}) & \textit{Usuario} (\texttt{id\_\allowbreak{}usuario}) & Sí \\ \hline
\textit{Codigo\-Segundo\-Factor} & (\texttt{id\_\allowbreak{}usuario}) & \textit{Usuario} (\texttt{id\_\allowbreak{}usuario}) & Sí \\ \hline
\textit{Token\-Verificacion\-Correo} & (\texttt{id\_\allowbreak{}usuario}) & \textit{Usuario} (\texttt{id\_\allowbreak{}usuario}) & Sí \\ \hline
\textit{Token\-Recuperacion} & (\texttt{id\_\allowbreak{}usuario}) & \textit{Usuario} (\texttt{id\_\allowbreak{}usuario}) & Sí \\ \hline
\textit{Aceptacion\-Terminos} & (\texttt{id\_\allowbreak{}usuario}) & \textit{Usuario} (\texttt{id\_\allowbreak{}usuario}) & Sí \\ \hline
\textit{Historial\-Nickname} & (\texttt{id\_\allowbreak{}usuario}) & \textit{Usuario} (\texttt{id\_\allowbreak{}usuario}) & Sí \\ \hline
\textit{Avatar} & (\texttt{id\_\allowbreak{}usuario}) & \textit{Usuario} (\texttt{id\_\allowbreak{}usuario}) & Sí \\ \hline
\textit{Enlace\-Red} & (\texttt{id\_\allowbreak{}usuario}) & \textit{Usuario} (\texttt{id\_\allowbreak{}usuario}) & Sí \\ \hline
\textit{Usuario\-Objeto} & (\texttt{id\_\allowbreak{}usuario}) & \textit{Usuario} (\texttt{id\_\allowbreak{}usuario}) & Sí \\ \hline
\textit{Usuario\-Objeto} & (\texttt{id\_\allowbreak{}objeto}) & \textit{Objeto\-Cosmetico} (\texttt{id\_\allowbreak{}objeto}) & Sí \\ \hline
\textit{Equipamiento} & (\texttt{id\_\allowbreak{}usuario}, \texttt{id\_\allowbreak{}objeto}) & \textit{Usuario\-Objeto} (\texttt{id\_\allowbreak{}usuario}, \texttt{id\_\allowbreak{}objeto}) & Sí \\ \hline
\textit{Equipamiento} & (\texttt{id\_\allowbreak{}objeto}, \texttt{id\_\allowbreak{}ranura}) & \textit{Objeto\-Cosmetico} (\texttt{id\_\allowbreak{}objeto}, \texttt{id\_\allowbreak{}ranura}) & Sí \\ \hline
\textit{Equipamiento} & (\texttt{id\_\allowbreak{}ranura}) & \textit{Ranura} (\texttt{id\_\allowbreak{}ranura}) & Sí \\ \hline
\textit{Partida} & (\texttt{id\_\allowbreak{}modo}) & \textit{Modo} (\texttt{id\_\allowbreak{}modo}) & Sí \\ \hline
\textit{Partida} & (\texttt{id\_\allowbreak{}nivel\_\allowbreak{}ia}) & \textit{Nivel\-IA} (\texttt{id\_\allowbreak{}nivel\_\allowbreak{}ia}) & No \\ \hline
\textit{Partida} & (\texttt{id\_\allowbreak{}partida}, \texttt{id\_\allowbreak{}ganador}) & \textit{Participacion} (\texttt{id\_\allowbreak{}partida}, \texttt{id\_\allowbreak{}usuario}) & No \\ \hline
\textit{Participacion} & (\texttt{id\_\allowbreak{}partida}) & \textit{Partida} (\texttt{id\_\allowbreak{}partida}) & Sí \\ \hline
\textit{Participacion} & (\texttt{id\_\allowbreak{}usuario}) & \textit{Usuario} (\texttt{id\_\allowbreak{}usuario}) & Sí \\ \hline
\textit{Participacion\-Objeto} & (\texttt{id\_\allowbreak{}partida}, \texttt{id\_\allowbreak{}usuario}) & \textit{Participacion} (\texttt{id\_\allowbreak{}partida}, \texttt{id\_\allowbreak{}usuario}) & Sí \\ \hline
\textit{Participacion\-Objeto} & (\texttt{id\_\allowbreak{}objeto}, \texttt{id\_\allowbreak{}ranura}) & \textit{Objeto\-Cosmetico} (\texttt{id\_\allowbreak{}objeto}, \texttt{id\_\allowbreak{}ranura}) & Sí \\ \hline
\textit{Jugada} & (\texttt{id\_\allowbreak{}partida}) & \textit{Partida} (\texttt{id\_\allowbreak{}partida}) & Sí \\ \hline
\textit{Jugada} & (\texttt{id\_\allowbreak{}partida}, \texttt{id\_\allowbreak{}usuario}) & \textit{Participacion} (\texttt{id\_\allowbreak{}partida}, \texttt{id\_\allowbreak{}usuario}) & No \\ \hline
\textit{Oferta\-Tablas} & (\texttt{id\_\allowbreak{}partida}, \texttt{id\_\allowbreak{}ofertante}) & \textit{Participacion} (\texttt{id\_\allowbreak{}partida}, \texttt{id\_\allowbreak{}usuario}) & Sí \\ \hline
\textit{Enlace\-Espectador} & (\texttt{id\_\allowbreak{}partida}, \texttt{id\_\allowbreak{}creador}) & \textit{Participacion} (\texttt{id\_\allowbreak{}partida}, \texttt{id\_\allowbreak{}usuario}) & Sí \\ \hline
\textit{Cola\-Emparejamiento} & (\texttt{id\_\allowbreak{}usuario}) & \textit{Usuario} (\texttt{id\_\allowbreak{}usuario}) & Sí \\ \hline
\textit{Cola\-Emparejamiento} & (\texttt{id\_\allowbreak{}modo}) & \textit{Modo} (\texttt{id\_\allowbreak{}modo}) & Sí \\ \hline
\textit{Sala} & (\texttt{id\_\allowbreak{}anfitrion}) & \textit{Usuario} (\texttt{id\_\allowbreak{}usuario}) & Sí \\ \hline
\textit{Sala} & (\texttt{id\_\allowbreak{}modo}) & \textit{Modo} (\texttt{id\_\allowbreak{}modo}) & Sí \\ \hline
\textit{Sala} & (\texttt{id\_\allowbreak{}partida}) & \textit{Partida} (\texttt{id\_\allowbreak{}partida}) & No \\ \hline
\textit{Sala\-Participante} & (\texttt{id\_\allowbreak{}sala}) & \textit{Sala} (\texttt{id\_\allowbreak{}sala}) & Sí \\ \hline
\textit{Sala\-Participante} & (\texttt{id\_\allowbreak{}usuario}) & \textit{Usuario} (\texttt{id\_\allowbreak{}usuario}) & Sí \\ \hline
\textit{Invitacion} & (\texttt{id\_\allowbreak{}sala}) & \textit{Sala} (\texttt{id\_\allowbreak{}sala}) & Sí \\ \hline
\textit{Invitacion} & (\texttt{id\_\allowbreak{}emisor}) & \textit{Usuario} (\texttt{id\_\allowbreak{}usuario}) & Sí \\ \hline
\textit{Invitacion} & (\texttt{id\_\allowbreak{}destinatario}) & \textit{Usuario} (\texttt{id\_\allowbreak{}usuario}) & Sí \\ \hline
\textit{Mensaje} & (\texttt{id\_\allowbreak{}autor}) & \textit{Usuario} (\texttt{id\_\allowbreak{}usuario}) & Sí \\ \hline
\textit{Mensaje} & (\texttt{id\_\allowbreak{}partida}) & \textit{Partida} (\texttt{id\_\allowbreak{}partida}) & No \\ \hline
\textit{Mensaje} & (\texttt{id\_\allowbreak{}sala}) & \textit{Sala} (\texttt{id\_\allowbreak{}sala}) & No \\ \hline
\textit{Estadistica\-Modo} & (\texttt{id\_\allowbreak{}usuario}) & \textit{Usuario} (\texttt{id\_\allowbreak{}usuario}) & Sí \\ \hline
\textit{Estadistica\-Modo} & (\texttt{id\_\allowbreak{}modo}) & \textit{Modo} (\texttt{id\_\allowbreak{}modo}) & Sí \\ \hline
\textit{Estadistica\-Modo} & (\texttt{id\_\allowbreak{}division}) & \textit{Division} (\texttt{id\_\allowbreak{}division}) & No \\ \hline
\textit{Historial\-Division} & (\texttt{id\_\allowbreak{}usuario}, \texttt{id\_\allowbreak{}modo}) & \textit{Estadistica\-Modo} (\texttt{id\_\allowbreak{}usuario}, \texttt{id\_\allowbreak{}modo}) & Sí \\ \hline
\textit{Historial\-Division} & (\texttt{id\_\allowbreak{}division\_\allowbreak{}anterior}) & \textit{Division} (\texttt{id\_\allowbreak{}division}) & No \\ \hline
\textit{Historial\-Division} & (\texttt{id\_\allowbreak{}division\_\allowbreak{}nueva}) & \textit{Division} (\texttt{id\_\allowbreak{}division}) & Sí \\ \hline
\textit{Caja} & (\texttt{id\_\allowbreak{}usuario}) & \textit{Usuario} (\texttt{id\_\allowbreak{}usuario}) & Sí \\ \hline
\textit{Caja} & (\texttt{id\_\allowbreak{}tipo\_\allowbreak{}caja}) & \textit{Tipo\-Caja} (\texttt{id\_\allowbreak{}tipo\_\allowbreak{}caja}) & Sí \\ \hline
\textit{Caja\-Contenido} & (\texttt{id\_\allowbreak{}caja}) & \textit{Caja} (\texttt{id\_\allowbreak{}caja}) & Sí \\ \hline
\textit{Caja\-Contenido} & (\texttt{id\_\allowbreak{}objeto}) & \textit{Objeto\-Cosmetico} (\texttt{id\_\allowbreak{}objeto}) & Sí \\ \hline
\textit{Movimiento\-Moneda} & (\texttt{id\_\allowbreak{}usuario}) & \textit{Usuario} (\texttt{id\_\allowbreak{}usuario}) & Sí \\ \hline
\textit{Movimiento\-Moneda} & (\texttt{id\_\allowbreak{}partida}) & \textit{Partida} (\texttt{id\_\allowbreak{}partida}) & No \\ \hline
\textit{Movimiento\-Moneda} & (\texttt{id\_\allowbreak{}objeto}) & \textit{Objeto\-Cosmetico} (\texttt{id\_\allowbreak{}objeto}) & No \\ \hline
\textit{Movimiento\-Moneda} & (\texttt{id\_\allowbreak{}caja}) & \textit{Caja} (\texttt{id\_\allowbreak{}caja}) & No \\ \hline
\textit{Amistad} & (\texttt{id\_\allowbreak{}usuario\_\allowbreak{}a}) & \textit{Usuario} (\texttt{id\_\allowbreak{}usuario}) & Sí \\ \hline
\textit{Amistad} & (\texttt{id\_\allowbreak{}usuario\_\allowbreak{}b}) & \textit{Usuario} (\texttt{id\_\allowbreak{}usuario}) & Sí \\ \hline
\textit{Solicitud} & (\texttt{id\_\allowbreak{}solicitante}) & \textit{Usuario} (\texttt{id\_\allowbreak{}usuario}) & Sí \\ \hline
\textit{Solicitud} & (\texttt{id\_\allowbreak{}destinatario}) & \textit{Usuario} (\texttt{id\_\allowbreak{}usuario}) & Sí \\ \hline
\textit{Silencio} & (\texttt{id\_\allowbreak{}usuario}) & \textit{Usuario} (\texttt{id\_\allowbreak{}usuario}) & Sí \\ \hline
\textit{Silencio} & (\texttt{id\_\allowbreak{}silenciado}) & \textit{Usuario} (\texttt{id\_\allowbreak{}usuario}) & Sí \\ \hline
\textit{Bloqueo} & (\texttt{id\_\allowbreak{}bloqueador}) & \textit{Usuario} (\texttt{id\_\allowbreak{}usuario}) & Sí \\ \hline
\textit{Bloqueo} & (\texttt{id\_\allowbreak{}bloqueado}) & \textit{Usuario} (\texttt{id\_\allowbreak{}usuario}) & Sí \\ \hline
\textit{Progreso\-Tutorial} & (\texttt{id\_\allowbreak{}usuario}) & \textit{Usuario} (\texttt{id\_\allowbreak{}usuario}) & Sí \\ \hline
\textit{Progreso\-Tutorial} & (\texttt{id\_\allowbreak{}leccion}) & \textit{Leccion} (\texttt{id\_\allowbreak{}leccion}) & Sí \\ \hline
\textit{Reporte} & (\texttt{id\_\allowbreak{}denunciante}) & \textit{Usuario} (\texttt{id\_\allowbreak{}usuario}) & Sí \\ \hline
\textit{Reporte} & (\texttt{id\_\allowbreak{}reportado}) & \textit{Usuario} (\texttt{id\_\allowbreak{}usuario}) & Sí \\ \hline
\textit{Reporte} & (\texttt{id\_\allowbreak{}motivo}) & \textit{Motivo\-Reporte} (\texttt{id\_\allowbreak{}motivo}) & Sí \\ \hline
\textit{Reporte} & (\texttt{id\_\allowbreak{}partida}) & \textit{Partida} (\texttt{id\_\allowbreak{}partida}) & No \\ \hline
\textit{Reporte} & (\texttt{id\_\allowbreak{}moderador}) & \textit{Usuario} (\texttt{id\_\allowbreak{}usuario}) & No \\ \hline
\textit{Sancion} & (\texttt{id\_\allowbreak{}usuario}) & \textit{Usuario} (\texttt{id\_\allowbreak{}usuario}) & Sí \\ \hline
\textit{Sancion} & (\texttt{id\_\allowbreak{}moderador}) & \textit{Usuario} (\texttt{id\_\allowbreak{}usuario}) & Sí \\ \hline
\textit{Sancion} & (\texttt{id\_\allowbreak{}reporte}) & \textit{Reporte} (\texttt{id\_\allowbreak{}reporte}) & No \\ \hline
\textit{Sancion} & (\texttt{id\_\allowbreak{}sancion\_\allowbreak{}sustituta}) & \textit{Sancion} (\texttt{id\_\allowbreak{}sancion}) & No \\ \hline
\textit{Apelacion} & (\texttt{id\_\allowbreak{}sancion}) & \textit{Sancion} (\texttt{id\_\allowbreak{}sancion}) & Sí \\ \hline
\textit{Apelacion} & (\texttt{id\_\allowbreak{}moderador}) & \textit{Usuario} (\texttt{id\_\allowbreak{}usuario}) & No \\ \hline
\textit{Bitacora\-Moderacion} & (\texttt{id\_\allowbreak{}moderador}) & \textit{Usuario} (\texttt{id\_\allowbreak{}usuario}) & Sí \\ \hline
\textit{Bitacora\-Moderacion} & (\texttt{id\_\allowbreak{}usuario\_\allowbreak{}afectado}) & \textit{Usuario} (\texttt{id\_\allowbreak{}usuario}) & No \\ \hline
\textit{Bitacora\-Acceso} & (\texttt{id\_\allowbreak{}usuario}) & \textit{Usuario} (\texttt{id\_\allowbreak{}usuario}) & No \\ \hline
\end{longtable}
\end{center}

